\documentclass{article}
\usepackage{amsmath, amssymb, amsthm, mathtools}

\setlength{\parindent}{0pt}

\newtheorem{claim}{Claim}

\begin{document}

\begin{claim}
    If $f : A \to B$ and $g : B \to C$ are bijections, then so is $g \circ f : A \to C$.
\end{claim}

\begin{proof}
    To show that $g \circ f$ is a bijection, we show that it is both injective and 
    surjective.

    \medskip
    \textbf{Injectivity.}
    Suppose $g(f(x_1)) = g(f(x_2))$. Since $g$ is injective, $f(x_1) = f(x_2)$. Since $f$
    is injective, $x_1 = x_2$. Hence $g \circ f$ is injective.

    \medskip
    \textbf{Surjectivity.}
    Let $c \in C$ be arbitrary. Since $g$ is surjective, there exists $b \in B$ such that
    $g(b) = c$. Since $f$ is surjective, there exists $a \in A$ such that $f(a) = b$. Then
    $(g \circ f)(a) = g(b) = c$. Hence $g \circ f$ is surjective.

    \medskip

    Since $g \circ f$ is both injective and surjective, it is a bijection.
\end{proof}

\begin{claim}
    If $f : A \to B$ is a bijection, then there exists a bijection $e : B \to A$ such that
    $e \circ f = I_A$.
\end{claim}

\begin{proof}
    Define $e : B \to A$ by $e(b) = a$, where $a \in A$ such that $f(a) = b$. Since $f$ is
    a bijection, for every $b \in B$ there exists a unique $a \in A$ such that $f(a) = b$.
    Hence $e$ is well-defined. To show that $e$ is a bijection, we show that it is both 
    injective and surjective.

    \medskip
    \textbf{Injectivity.}
    Suppose $e(b_1) = e(b_2)$. Let $a_1, a_2 \in A$ be such that $f(a_1) = b_1$ and $f(a_2) 
    = b_2$. Then by definition of $e$, $a_1 = a_2$. Since $f(a_1) = b_1$ and $f(a_2) = b_2$, 
    $b_1 = b_2$. Hence $e$ is injective.

    \medskip
    \textbf{Surjectivity.}
    Let $a \in A$ be arbitrary. By definition of $f$, $f(a) = b \in B$, so $e(b) = e(f(a)) 
    = a$. Hence $e$ is surjective.

    \medskip

    Therefore for any $a \in A$, $(e \circ f)(a) = e(f(a)) = a = I_A(a)$ by definition of $e$. 
    Thus such a bijection $e : B \to A$ exists. 
\end{proof}

\begin{claim}
    Graph isomorphism is an equivalence.
\end{claim}

\begin{proof}
    
\end{proof}

\end{document}
